\begin{lstlisting}[language=Ruby, caption={Vagrantfile configuration for Scenario 2 (PANOPTES Architecture)}, label={lst:vagrant-scenario2}]
# -*- mode: ruby -*-
# vi: set ft=ruby :

Vagrant.configure("2") do |config|
  config.vm.box = "ubuntu/jammy64"

  # ==========================================
  # 1. CONTROL PLANE (MASTER NODE) - PANOPTES
  # ==========================================
  config.vm.define "master" do |master|
    master.vm.hostname = "k8s-master-panoptes"
    master.vm.network "private_network", ip: "192.168.56.10"
    
    # Expose Grafana to the host machine for easy access
    master.vm.network "forwarded_port", guest: 80, host: 8080, auto_correct: true

    master.vm.provider "virtualbox" do |vb|
      vb.name = "k8s-master-panoptes"
      vb.memory = "4096"
      vb.cpus = 2
      vb.customize ["modifyvm", :id, "--groups", "/k8s-panoptes"]
    end

    # Base Provisioning and K3s Server
    master.vm.provision "shell", inline: <<-SHELL
      export DEBIAN_FRONTEND=noninteractive
      apt-get update -y
      
      # Install OS dependencies and network tools
      echo "[INFO] Installing base dependencies (curl, jq, git)..."
      apt-get install -y software-properties-common curl jq git python3-pip \
                         linux-headers-$(uname -r) linux-tools-$(uname -r) bpfcc-tools

      # Install Ansible natively on the Master node
      echo "[INFO] Installing Ansible and Kubernetes modules..."
      add-apt-repository --yes --update ppa:ansible/ansible
      apt-get install -y ansible
      ansible-galaxy collection install kubernetes.core
      
      echo "[INFO] Installing Python dependencies for Ansible..."
      pip3 install kubernetes PyYAML jsonpatch

      # Remove conflicting snap packages and install official Helm 3
      snap remove helm 2>/dev/null || true
      echo "[INFO] Installing official Helm 3..."
      curl -fsSL -o get_helm.sh https://raw.githubusercontent.com/helm/helm/main/scripts/get-helm-3
      bash get_helm.sh

      # Discover VirtualBox private network interface
      FLANNEL_IFACE=$(ip -4 a | grep 192.168.56 | awk '{print $NF}')
      
      echo "[INFO] Installing K3s Server..."
      curl -sfL https://get.k3s.io | INSTALL_K3S_EXEC="--write-kubeconfig-mode 644 --node-ip 192.168.56.10 --disable traefik --flannel-iface=$FLANNEL_IFACE" sh -

      # Configure kubeconfig for root and vagrant users
      mkdir -p /root/.kube /home/vagrant/.kube
      cp /etc/rancher/k3s/k3s.yaml /root/.kube/config
      cp /etc/rancher/k3s/k3s.yaml /home/vagrant/.kube/config
      chown -R vagrant:vagrant /home/vagrant/.kube

      # Share K3s token for workers to join
      sleep 15
      cp /var/lib/rancher/k3s/server/node-token /vagrant/node-token
      echo "[INFO] Master provisioned successfully. Token generated."
    SHELL
  end

  # ==========================================
  # 2. DATA PLANE (WORKER NODES)
  # ==========================================
  (1..2).each do |i|
    config.vm.define "worker#{i}" do |worker|
      worker.vm.hostname = "k8s-worker#{i}-panoptes"
      worker.vm.network "private_network", ip: "192.168.56.1#{i}"

      worker.vm.provider "virtualbox" do |vb|
        vb.name = "k8s-worker#{i}-panoptes"
        vb.memory = "4096"
        vb.cpus = 2
        vb.customize ["modifyvm", :id, "--groups", "/k8s-panoptes"]
      end

      worker.vm.provision "shell", inline: <<-SHELL
        export DEBIAN_FRONTEND=noninteractive
        apt-get update -y

        echo "[INFO] Installing base dependencies on Worker..."
        apt-get install -y curl linux-headers-$(uname -r) linux-tools-$(uname -r) bpfcc-tools

        FLANNEL_IFACE=$(ip -4 a | grep 192.168.56 | awk '{print $NF}')
        
        echo "[INFO] Waiting for the Master Node to generate the join token..."
        while [ ! -f /vagrant/node-token ]; do sleep 2; done
        TOKEN=$(cat /vagrant/node-token)

        echo "[INFO] Joining the K3s cluster as an Agent..."
        curl -sfL https://get.k3s.io | K3S_URL=https://192.168.56.10:6443 K3S_TOKEN=$TOKEN INSTALL_K3S_EXEC="--node-ip 192.168.56.1#{i} --flannel-iface=$FLANNEL_IFACE" sh -
      SHELL
    end
  end
end
\end{lstlisting}